\documentclass[10pt,a4paper]{amsart}
\usepackage{amsmath,amssymb,amsthm,graphicx}
% For dealing with references we use the comment environment
\usepackage{verbatim}
\newenvironment{reference}{\comment}{\endcomment}
%\newenvironment{reference}{}{}
\newenvironment{slogan}{\comment}{\endcomment}
\newenvironment{history}{\comment}{\endcomment}
% For commutative diagrams we use Xy-pic
\usepackage[all]{xy}

% We use 2cell for 2-commutative diagrams.
\xyoption{2cell}
\UseAllTwocells
\theoremstyle{plain}
\newtheorem{theorem}[subsection]{Theorem}
\newtheorem{proposition}[subsection]{Proposition}
\newtheorem{lemma}[subsection]{Lemma}

\theoremstyle{definition}
\newtheorem{definition}[subsection]{Definition}
\newtheorem{example}[subsection]{Example}
\newtheorem{exercise}[subsection]{Exercise}
\newtheorem{situation}[subsection]{Situation}

\theoremstyle{remark}
\newtheorem{remark}[subsection]{Remark}
\newtheorem{remarks}[subsection]{Remarks}

\numberwithin{equation}{subsection}
% Macros
%
\def\lim{\mathop{\mathrm{lim}}\nolimits}
\def\colim{\mathop{\mathrm{colim}}\nolimits}
\def\Spec{\mathop{\mathrm{Spec}}}
\def\Hom{\mathop{\mathrm{Hom}}\nolimits}
\def\Ext{\mathop{\mathrm{Ext}}\nolimits}
\def\SheafHom{\mathop{\mathcal{H}\!\mathit{om}}\nolimits}
\def\SheafExt{\mathop{\mathcal{E}\!\mathit{xt}}\nolimits}
\def\Sch{\mathit{Sch}}
\def\Mor{\mathop{\mathrm{Mor}}\nolimits}
\def\Ob{\mathop{\mathrm{Ob}}\nolimits}
\def\Sh{\mathop{\mathit{Sh}}\nolimits}
\def\NL{\mathop{N\!L}\nolimits}
\def\CH{\mathop{\mathrm{CH}}\nolimits}
\def\proetale{{pro\text{-}\acute{e}tale}}
\def\etale{{\acute{e}tale}}
\def\QCoh{\mathit{QCoh}}
\def\Ker{\mathop{\mathrm{Ker}}}
\def\Im{\mathop{\mathrm{Im}}}
\def\Coker{\mathop{\mathrm{Coker}}}
\def\Coim{\mathop{\mathrm{Coim}}}

% Boxtimes
%
\DeclareMathSymbol{\boxtimes}{\mathbin}{AMSa}{"02}

%
% Macros for moduli stacks/spaces
%
\def\QCohstack{\mathcal{QC}\!\mathit{oh}}
\def\Cohstack{\mathcal{C}\!\mathit{oh}}
\def\Spacesstack{\mathcal{S}\!\mathit{paces}}
\def\Quotfunctor{\mathrm{Quot}}
\def\Hilbfunctor{\mathrm{Hilb}}
\def\Curvesstack{\mathcal{C}\!\mathit{urves}}
\def\Polarizedstack{\mathcal{P}\!\mathit{olarized}}
\def\Complexesstack{\mathcal{C}\!\mathit{omplexes}}
% \Pic is the operator that assigns to X its picard group, usage \Pic(X)
% \Picardstack_{X/B} denotes the Picard stack of X over B
% \Picardfunctor_{X/B} denotes the Picard functor of X over B
\def\Pic{\mathop{\mathrm{Pic}}\nolimits}
\def\Picardstack{\mathcal{P}\!\mathit{ic}}
\def\Picardfunctor{\mathrm{Pic}}
\def\Deformationcategory{\mathcal{D}\!\mathit{ef}}
\usepackage{cleveref}
\crefname{originalreference}{original reference}{original references}
\Crefname{originalreference}{Original reference}{Original references}
\makeatletter
\newlabel{situation-existence}{{08B7}{source}{}{sourcepointer-1}{}}
\newlabel{situation-existence@cref}{{[originalreference][0][]08B7}{[1][0][]source}{}{}{}}
\makeatother
\crefname{originalequation}{original equation}{original equations}
\Crefname{originalequation}{Original equation}{Original equations}
\makeatletter
\newlabel{equation-completion-functor-proper-over-A}{{08B8}{source}{}{sourcepointer-2}{}}
\newlabel{equation-completion-functor-proper-over-A@cref}{{[originalequation][0][]08B8}{[1][0][]source}{}{}{}}
\makeatother
\crefname{originalreference}{original reference}{original references}
\Crefname{originalreference}{Original reference}{Original references}
\makeatletter
\newlabel{spaces-cohomology-lemma-i-star-equivalence}{{08AM}{source}{}{sourcepointer-3}{}}
\newlabel{spaces-cohomology-lemma-i-star-equivalence@cref}{{[originalreference][0][]08AM}{[1][0][]source}{}{}{}}
\makeatother
\crefname{originallemma}{original lemma}{original lemmas}
\Crefname{originallemma}{Original lemma}{Original lemmas}
\makeatletter
\newlabel{lemma-fully-faithful}{{08B9}{source}{}{sourcepointer-4}{}}
\newlabel{lemma-fully-faithful@cref}{{[originallemma][0][]08B9}{[1][0][]source}{}{}{}}
\makeatother
\crefname{originalreference}{original reference}{original references}
\Crefname{originalreference}{Original reference}{Original references}
\makeatletter
\newlabel{spaces-cohomology-lemma-acc-coherent}{{07UJ}{source}{}{sourcepointer-5}{}}
\newlabel{spaces-cohomology-lemma-acc-coherent@cref}{{[originalreference][0][]07UJ}{[1][0][]source}{}{}{}}
\makeatother
\crefname{originallemma}{original lemma}{original lemmas}
\Crefname{originallemma}{Original lemma}{Original lemmas}
\makeatletter
\newlabel{lemma-chow-noetherian-separated}{{088U}{source}{}{sourcepointer-6}{}}
\newlabel{lemma-chow-noetherian-separated@cref}{{[originallemma][0][]088U}{[1][0][]source}{}{}{}}
\makeatother
\crefname{originallemma}{original lemma}{original lemmas}
\Crefname{originallemma}{Original lemma}{Original lemmas}
\makeatletter
\newlabel{morphisms-lemma-H-quasi-projective-open-H-projective}{{01WA}{source}{}{sourcepointer-7}{}}
\newlabel{morphisms-lemma-H-quasi-projective-open-H-projective@cref}{{[originallemma][0][]01WA}{[1][0][]source}{}{}{}}
\makeatother
\crefname{originalreference}{original reference}{original references}
\Crefname{originalreference}{Original reference}{Original references}
\makeatletter
\newlabel{spaces-morphisms-lemma-image-proper-is-proper}{{08AJ}{source}{}{sourcepointer-8}{}}
\newlabel{spaces-morphisms-lemma-image-proper-is-proper@cref}{{[originalreference][0][]08AJ}{[1][0][]source}{}{}{}}
\makeatother
\crefname{originalreference}{original reference}{original references}
\Crefname{originalreference}{Original reference}{Original references}
\makeatletter
\newlabel{spaces-morphisms-lemma-composition-proper}{{04XY}{source}{}{sourcepointer-9}{}}
\newlabel{spaces-morphisms-lemma-composition-proper@cref}{{[originalreference][0][]04XY}{[1][0][]source}{}{}{}}
\makeatother
\crefname{originallemma}{original lemma}{original lemmas}
\Crefname{originallemma}{Original lemma}{Original lemmas}
\makeatletter
\newlabel{morphisms-lemma-image-proper-scheme-closed}{{01W6}{source}{}{sourcepointer-10}{}}
\newlabel{morphisms-lemma-image-proper-scheme-closed@cref}{{[originallemma][0][]01W6}{[1][0][]source}{}{}{}}
\makeatother
\crefname{originallemma}{original lemma}{original lemmas}
\Crefname{originallemma}{Original lemma}{Original lemmas}
\makeatletter
\newlabel{coherent-lemma-existence-projective}{{0885}{source}{}{sourcepointer-11}{}}
\newlabel{coherent-lemma-existence-projective@cref}{{[originallemma][0][]0885}{[1][0][]source}{}{}{}}
\makeatother
\crefname{originallemma}{original lemma}{original lemmas}
\Crefname{originallemma}{Original lemma}{Original lemmas}
\makeatletter
\newlabel{algebra-lemma-radical-completion}{{05GI}{source}{}{sourcepointer-12}{}}
\newlabel{algebra-lemma-radical-completion@cref}{{[originallemma][0][]05GI}{[1][0][]source}{}{}{}}
\makeatother
\crefname{originalreference}{original reference}{original references}
\Crefname{originalreference}{Original reference}{Original references}
\makeatletter
\newlabel{spaces-cohomology-lemma-proper-pushforward-coherent}{{08AR}{source}{}{sourcepointer-13}{}}
\newlabel{spaces-cohomology-lemma-proper-pushforward-coherent@cref}{{[originalreference][0][]08AR}{[1][0][]source}{}{}{}}
\makeatother
\crefname{originallemma}{original lemma}{original lemmas}
\Crefname{originallemma}{Original lemma}{Original lemmas}
\makeatletter
\newlabel{lemma-inverse-systems-push-pull}{{08BD}{source}{}{sourcepointer-14}{}}
\newlabel{lemma-inverse-systems-push-pull@cref}{{[originallemma][0][]08BD}{[1][0][]source}{}{}{}}
\makeatother
\crefname{originallemma}{original lemma}{original lemmas}
\Crefname{originallemma}{Original lemma}{Original lemmas}
\makeatletter
\newlabel{lemma-existence-tricky}{{08BC}{source}{}{sourcepointer-15}{}}
\newlabel{lemma-existence-tricky@cref}{{[originallemma][0][]08BC}{[1][0][]source}{}{}{}}
\makeatother
\title{theorem grothendieck existence}
\author{The Stacks Project Authors}
\date{stacks-project:a04446e57ec1fbc252a871afcec7752fb2807b14; GNU Free Documentation License 1.2 (as published in the Stacks Project COPYING)}
\makeatletter
\newlabel{spaces-perfect-lemma-module-support-proper-over-base}{{0CZI}{source}{}{context-source-1}{}}
\makeatother
\makeatletter
\newlabel{equation-completion-functor}{{08B4}{source}{}{context-source-2}{}}
\makeatother
\begin{document}
\maketitle
\noindent\textit{Editorial note.} All mathematical statement and proof text is reproduced verbatim from the author's chapter spaces-more-morphisms.tex of the Stacks Project, fixed at commit a04446e57ec1fbc252a871afcec7752fb2807b14. The pinned preamble is supplied as its own file; the document wrapper class, the xr-hyper wiring and the \textbackslash{}externaldocument list are omitted because the assembler supplies the class.\par\medskip
\noindent\textit{Editorial note.} Numbering inside this standalone excerpt is editorial, not the full Stacks book print numbering.\par\medskip
\noindent\textit{Editorial note.} The candidate was previously judged H2 by the initial-corpus stage; that judgement and its reason are reused verbatim and were not re-derived here. Original candidate key: S001-spaces-more-morphisms-theorem-grothendieck-existence.\par\medskip
\noindent\textit{Editorial note.} Reference labels in the proof resolve to navigation pointers to the original author source; these pointers do not reproduce proofs of the prerequisites.\par\medskip
\noindent\textit{Allowed standard prerequisites (editorial).}\begin{itemize}
\item Prerequisites named by the author inside the counted span. The counted span is reproduced verbatim, including the author's own citations. Its prerequisites are the earlier results of the same fixed Stacks Project version that the author cites there and in the chapter's lead-in.
\end{itemize}\medskip
\ifdefined\hypertarget\hypertarget{sourcepointer-1}{}\fi
\noindent\textit{Source reference.} Original reference 08B7: sources/spaces-more-morphisms.tex, source lines 9682--9682. Navigation pointer used in the counted span; the span contains the pointer only and the cited result is not reproduced inside it. Registered as a non-load-bearing navigation stub in the established form for Stacks material. No bibliography citation is involved.\par\medskip
\ifdefined\hypertarget\hypertarget{sourcepointer-2}{}\fi
\noindent\textit{Source reference.} Original equation 08B8: sources/spaces-more-morphisms.tex, source lines 9710--9710. Navigation pointer used in the counted span; the span contains the pointer only and the cited result is not reproduced inside it. Registered as a non-load-bearing navigation stub in the established form for Stacks material. No bibliography citation is involved.\par\medskip
\ifdefined\hypertarget\hypertarget{sourcepointer-3}{}\fi
\noindent\textit{Source reference.} Original reference 08AM: sources/spaces-cohomology.tex, source lines 2042--2042. Navigation pointer used in the counted span; the span contains the pointer only and the cited result is not reproduced inside it. Registered as a non-load-bearing navigation stub in the established form for Stacks material. No bibliography citation is involved.\par\medskip
\ifdefined\hypertarget\hypertarget{sourcepointer-4}{}\fi
\noindent\textit{Source reference.} Original lemma 08B9: sources/spaces-more-morphisms.tex, source lines 9718--9718. Navigation pointer used in the counted span; the span contains the pointer only and the cited result is not reproduced inside it. Registered as a non-load-bearing navigation stub in the established form for Stacks material. No bibliography citation is involved.\par\medskip
\ifdefined\hypertarget\hypertarget{sourcepointer-5}{}\fi
\noindent\textit{Source reference.} Original reference 07UJ: sources/spaces-cohomology.tex, source lines 2096--2096. Navigation pointer used in the counted span; the span contains the pointer only and the cited result is not reproduced inside it. Registered as a non-load-bearing navigation stub in the established form for Stacks material. No bibliography citation is involved.\par\medskip
\ifdefined\hypertarget\hypertarget{sourcepointer-6}{}\fi
\noindent\textit{Source reference.} Original lemma 088U: sources/spaces-more-morphisms.tex, source lines 9275--9275. Navigation pointer used in the counted span; the span contains the pointer only and the cited result is not reproduced inside it. Registered as a non-load-bearing navigation stub in the established form for Stacks material. No bibliography citation is involved.\par\medskip
\ifdefined\hypertarget\hypertarget{sourcepointer-7}{}\fi
\noindent\textit{Source reference.} Original lemma 01WA: sources/morphisms.tex, source lines 10703--10703. Navigation pointer used in the counted span; the span contains the pointer only and the cited result is not reproduced inside it. Registered as a non-load-bearing navigation stub in the established form for Stacks material. No bibliography citation is involved.\par\medskip
\ifdefined\hypertarget\hypertarget{sourcepointer-8}{}\fi
\noindent\textit{Source reference.} Original reference 08AJ: sources/spaces-morphisms.tex, source lines 7994--7994. Navigation pointer used in the counted span; the span contains the pointer only and the cited result is not reproduced inside it. Registered as a non-load-bearing navigation stub in the established form for Stacks material. No bibliography citation is involved.\par\medskip
\ifdefined\hypertarget\hypertarget{sourcepointer-9}{}\fi
\noindent\textit{Source reference.} Original reference 04XY: sources/spaces-morphisms.tex, source lines 7928--7928. Navigation pointer used in the counted span; the span contains the pointer only and the cited result is not reproduced inside it. Registered as a non-load-bearing navigation stub in the established form for Stacks material. No bibliography citation is involved.\par\medskip
\ifdefined\hypertarget\hypertarget{sourcepointer-10}{}\fi
\noindent\textit{Source reference.} Original lemma 01W6: sources/morphisms.tex, source lines 9949--9949. Navigation pointer used in the counted span; the span contains the pointer only and the cited result is not reproduced inside it. Registered as a non-load-bearing navigation stub in the established form for Stacks material. No bibliography citation is involved.\par\medskip
\ifdefined\hypertarget\hypertarget{sourcepointer-11}{}\fi
\noindent\textit{Source reference.} Original lemma 0885: sources/coherent.tex, source lines 6812--6812. Navigation pointer used in the counted span; the span contains the pointer only and the cited result is not reproduced inside it. Registered as a non-load-bearing navigation stub in the established form for Stacks material. No bibliography citation is involved.\par\medskip
\ifdefined\hypertarget\hypertarget{sourcepointer-12}{}\fi
\noindent\textit{Source reference.} Original lemma 05GI: sources/algebra.tex, source lines 22734--22734. Navigation pointer used in the counted span; the span contains the pointer only and the cited result is not reproduced inside it. Registered as a non-load-bearing navigation stub in the established form for Stacks material. No bibliography citation is involved.\par\medskip
\ifdefined\hypertarget\hypertarget{sourcepointer-13}{}\fi
\noindent\textit{Source reference.} Original reference 08AR: sources/spaces-cohomology.tex, source lines 3618--3618. Navigation pointer used in the counted span; the span contains the pointer only and the cited result is not reproduced inside it. Registered as a non-load-bearing navigation stub in the established form for Stacks material. No bibliography citation is involved.\par\medskip
\ifdefined\hypertarget\hypertarget{sourcepointer-14}{}\fi
\noindent\textit{Source reference.} Original lemma 08BD: sources/spaces-more-morphisms.tex, source lines 9941--9941. Navigation pointer used in the counted span; the span contains the pointer only and the cited result is not reproduced inside it. Registered as a non-load-bearing navigation stub in the established form for Stacks material. No bibliography citation is involved.\par\medskip
\ifdefined\hypertarget\hypertarget{sourcepointer-15}{}\fi
\noindent\textit{Source reference.} Original lemma 08BC: sources/spaces-more-morphisms.tex, source lines 9828--9828. Navigation pointer used in the counted span; the span contains the pointer only and the cited result is not reproduced inside it. Registered as a non-load-bearing navigation stub in the established form for Stacks material. No bibliography citation is involved.\par\medskip
\noindent\textit{Reading revision r002.} Original author context is supplied below and editorial reference bindings are qualified by the chapter. Mathematical statement and proof are unchanged. Numbering is editorial; cited inputs remain in the complete fixed source.\par\medskip
\subsection*{Original author context (uncounted)}
\begin{situation}
\label{context-situation-existence}
Here $A$ is a Noetherian ring complete with respect to an ideal $I$.
Also $f : X \to \Spec(A)$ is a finite type separated morphism of
algebraic spaces and $\mathcal{I} = I\mathcal{O}_X$.
\end{situation}

\noindent
In this situation we denote
$$
\textit{Coh}_{\text{support proper over } A}(\mathcal{O}_X)
$$
be the full subcategory of $\textit{Coh}(\mathcal{O}_X)$
consisting of those coherent $\mathcal{O}_X$-modules whose
support is proper over $\Spec(A)$, or equivalently whose
scheme theoretic support is proper over $\Spec(A)$, see
Derived Categories of Spaces, Lemma
\ref{spaces-perfect-lemma-module-support-proper-over-base}.
Similarly, we let
$$
\textit{Coh}_{\text{support proper over } A}(X, \mathcal{I})
$$
be the full subcategory of $\textit{Coh}(X, \mathcal{I})$
consisting of those objects $(\mathcal{F}_n)$ such that
the support of $\mathcal{F}_1$ is proper over $\Spec(A)$.
Since the support of a quotient module is contained in the support
of the module, it follows that (\ref{equation-completion-functor})
induces a functor
\begin{equation}
\label{context-equation-completion-functor-proper-over-A}
\textit{Coh}_{\text{support proper over }A}(\mathcal{O}_X)
\longrightarrow
\textit{Coh}_{\text{support proper over }A}(X, \mathcal{I})
\end{equation}

\subsection*{Original author statement and proof}
\begin{theorem}[Grothendieck's existence theorem]
\label{theorem-grothendieck-existence}
In Situation \ref{situation-existence} the functor
(\ref{equation-completion-functor-proper-over-A})
is an equivalence.
\end{theorem}

\begin{proof}
We will use the equivalence of categories of
Cohomology of Spaces, Lemma \ref{spaces-cohomology-lemma-i-star-equivalence}
without further mention in the proof of the theorem.
By Lemma \ref{lemma-fully-faithful} the functor is fully faithful.
Thus we need to prove the functor is essentially surjective.

\medskip\noindent
Consider the collection $\Xi$ of quasi-coherent sheaves of ideals
$\mathcal{K} \subset \mathcal{O}_X$ such that the statement holds
for every object $(\mathcal{F}_n)$ of
$\textit{Coh}_{\text{support proper over }A}(X, \mathcal{I})$
annihilated by $\mathcal{K}$. We want to show $(0)$ is in $\Xi$.
If not, then since $X$ is Noetherian there exists a maximal
quasi-coherent sheaf of ideals $\mathcal{K}$ not in $\Xi$, see
Cohomology of Spaces, Lemma \ref{spaces-cohomology-lemma-acc-coherent}.
After replacing $X$ by the closed subscheme of $X$
corresponding to $\mathcal{K}$ we may assume that every nonzero
$\mathcal{K}$ is in $\Xi$. Let $(\mathcal{F}_n)$ be an object of
$\textit{Coh}_{\text{support proper over }A}(X, \mathcal{I})$.
We will show that this object is in the essential image, thereby
completing the proof of the theorem.

\medskip\noindent
Apply Chow's lemma (Lemma \ref{lemma-chow-noetherian-separated})
to find a proper surjective morphism $f : Y \to X$ which is an isomorphism
over a dense open $U \subset X$ such that $Y$ is H-quasi-projective
over $A$. Note that $Y$ is a scheme and $f$ representable.
Choose an open immersion $j : Y \to Y'$ with $Y'$ projective over $A$, see
Morphisms, Lemma \ref{morphisms-lemma-H-quasi-projective-open-H-projective}.
Let $T_n$ be the scheme theoretic support of $\mathcal{F}_n$.
Note that $|T_n| = |T_1|$, hence $T_n$ is proper over $A$ for all $n$
(Morphisms of Spaces, Lemma
\ref{spaces-morphisms-lemma-image-proper-is-proper}).
Then $f^*\mathcal{F}_n$ is supported on the closed subscheme
$f^{-1}T_n$ which is proper over $A$ (by
Morphisms of Spaces, Lemma \ref{spaces-morphisms-lemma-composition-proper}
and properness of $f$).
In particular, the composition $f^{-1}T_n \to Y \to Y'$ is closed
(Morphisms, Lemma \ref{morphisms-lemma-image-proper-scheme-closed}).
Let $T'_n \subset Y'$ be the corresponding closed subscheme;
it is contained in the open subscheme $Y$ and equal to $f^{-1}T_n$
as a closed subscheme of $Y$. Let $\mathcal{F}_n'$
be the coherent $\mathcal{O}_{Y'}$-module corresponding to
$f^*\mathcal{F}_n$ viewed as a coherent module on $Y'$ via
the closed immersion $f^{-1}T_n = T'_n \subset Y'$.
Then $(\mathcal{F}_n')$
is an object of $\textit{Coh}(Y', I\mathcal{O}_{Y'})$.
By the projective case of Grothendieck's existence theorem
(Cohomology of Schemes, Lemma \ref{coherent-lemma-existence-projective})
there exists a coherent $\mathcal{O}_{Y'}$-module
$\mathcal{F}'$ and an isomorphism
$(\mathcal{F}')^\wedge \cong (\mathcal{F}'_n)$ in
$\textit{Coh}(Y', I\mathcal{O}_{Y'})$.
Let $Z' \subset Y'$ be the scheme theoretic support of $\mathcal{F}'$.
Since $\mathcal{F}'/I\mathcal{F}' = \mathcal{F}'_1$ we see
that $Z' \cap V(I\mathcal{O}_{Y'}) = T'_1$ set-theoretically.
The structure morphism $p' : Y' \to \Spec(A)$ is proper, hence
$p'(Z' \cap (Y' \setminus Y))$ is closed in $\Spec(A)$.
If nonempty, then it would contain a point of $V(I)$
as $I$ is contained in the Jacobson radical of $A$
(Algebra, Lemma \ref{algebra-lemma-radical-completion}).
But we've seen above that $Z' \cap (p')^{-1}V(I) = T'_1 \subset Y$
hence we conclude that $Z' \subset Y$. Thus $\mathcal{F}'|_Y$
is supported on a closed subscheme of $Y$ proper over $A$.

\medskip\noindent
Let $\mathcal{K}$ be the quasi-coherent sheaf of ideals cutting
out the reduced complement $X \setminus U$.
By Cohomology of Spaces, Lemma
\ref{spaces-cohomology-lemma-proper-pushforward-coherent}
the $\mathcal{O}_X$-module $\mathcal{H} = f_*\mathcal{F}'$ is coherent
and by Lemma \ref{lemma-inverse-systems-push-pull}
there exists a morphism $\alpha : (\mathcal{F}_n) \to \mathcal{H}^\wedge$
in the category $\textit{Coh}_{\text{support proper over } A}(X, \mathcal{I})$
whose kernel and cokernel are annihilated by a power of $\mathcal{K}$.
Let $Z_0 \subset X$ be the scheme theoretic support of $\mathcal{H}$.
It is clear that $|Z_0| \subset f(|Z'|)$. Hence
$Z_0 \to \Spec(A)$ is proper
(Morphisms of Spaces, Lemma
\ref{spaces-morphisms-lemma-image-proper-is-proper}).
Thus $\mathcal{H}$ is an object of 
$\textit{Coh}_{\text{support proper over } A}(\mathcal{O}_X)$.
Since each of the sheaves of ideals $\mathcal{K}^e$ is an element of
$\Xi$ we see that the assumptions of Lemma \ref{lemma-existence-tricky}
are satisfied and we conclude.
\end{proof}

\end{document}
